\clearpage
\phantomsection
\addcontentsline{toc}{chapter}{Simulado I: Revisão Integrada da Unidade I (AV1)}
\chapter*{Simulado e Revisão Integrada da Unidade I (Preparatório AV1: 30\%)}
\markboth{SIMULADO, UNIDADE I}{SIMULADO, UNIDADE I}

\begin{capObjetivos}
Este simulado integra os conteúdos dos \textbf{Capítulos 1 a 6} (\textbf{Unidade I, Território e Ecossistemas}), servindo como preparação direta para o Seminário e Fichamento Crítico da \textbf{AV1 (30\%)}. Resolva as \textbf{8 questões objetivas (A--E)} e as \textbf{2 questões dissertativas críticas} antes de consultar o gabarito comentado, o espelho dissertativo e a tabela de diagnóstico.
\end{capObjetivos}

\etapa{1}{Questões Objetivas Integradas (Capítulos 1 a 6)}

\textbf{Questão 1 (Cap. 1, Schumpeter, Dosi e Espaço Territorial)}
Considerando a articulação entre inovação tecnológica e desenvolvimento territorial (Schumpeter, Kuhn, Dosi e Milton Santos), assinale a alternativa \textbf{correta}:
\begin{alternativas}
  \item A destruição criadora schumpeteriana distribui seus benefícios de modo uniforme por todo o espaço geográfico, independentemente da densidade técnica local.
  \item Paradigmas e trajetórias tecnológicas (Dosi) orientam a busca por soluções técnicas e tendem a se concentrar em territórios com maior densidade de infraestrutura, conhecimento e normas (meio técnico-científico-informacional).
  \item A inovação incremental corresponde a rupturas radicais de paradigma científico que inviabilizam todas as firmas incumbentes de imediato.
  \item Na leitura de Milton Santos, o território é apenas um suporte físico passivo, sem influência sobre a apropriação do valor tecnológico.
  \item O empreendedorismo de base tecnológica independe de trajetórias tecnológicas prévias ou de instituições regionais de ensino e pesquisa.
\end{alternativas}

\textbf{Questão 2 (Cap. 2, Estruturas de Mercado e Eficiência Coletiva)}
Em um Arranjo Produtivo Local (APL) de confecções ou calçados formado por microempresas, a superação de desvantagens de escala diante de grandes concorrentes externos ocorre principalmente por meio da \textbf{eficiência coletiva}, que combina:
\begin{alternativas}
  \item Monopólio legal concedido pelo município e eliminação de fornecedores especializados locais.
  \item Externalidades passivas da aglomeração (como concentração de mão de obra e fornecedores) somadas a ações conjuntas deliberadas (compras coletivas, consórcios de exportação e laboratórios compartilhados).
  \item Guerra fiscal predatória entre as próprias empresas do arranjo para reduzir salários abaixo do mínimo legal.
  \item Integração vertical absoluta de cada microempresa, proibindo a subcontratação local.
  \item Substituição integral da cooperação territorial pela importação de produtos acabados.
\end{alternativas}

\textbf{Questão 3 (Cap. 3, Ecossistemas Empreendedores e Proximidades de Boschma)}
Segundo Ron Boschma, a proximidade geográfica entre empresas e universidades \textbf{não garante} por si só o aprendizado inovativo, pois a cooperação eficaz exige equilíbrio entre outras dimensões de proximidade. O risco de \textit{lock-in} (aprisionamento cognitivo e rotinas obsoletas) em um território ocorre tipicamente quando há:
\begin{alternativas}
  \item Excesso de proximidade cognitiva, organizacional e institucional, reduzindo a entrada de novas ideias e o questionamento de práticas antigas.
  \item Baixa distância geográfica combinada com alta diversidade de atores em Hélice Quíntupla.
  \item Existência de intermediários tecnológicos (NITs) conectando a região a redes internacionais de pesquisa.
  \item Presença simultânea de \textit{startups}, cooperativas e fundos de fomento público.
  \item Renovação periódica das trajetórias tecnológicas por meio de políticas orientadas por missões.
\end{alternativas}

\textbf{Questão 4 (Cap. 4, Governança em Hélice Quíntupla e Marco Legal de CT\&I)}
Uma empresa tributada pelo \textbf{Lucro Real} investiu $\reais{1.500.000,00}$ em P\&D elegível e obteve direito à exclusão adicional de $\alpha_{excl} = 80\%$ pela Lei do Bem (Lei n.\textsuperscript{o} 11.196/2005). Considerando alíquota efetiva combinada de IRPJ e CSLL de $34\%$, a economia fiscal incentivada ($E_{LeiBem}$) e o instrumento adequado caso um órgão público queira contratar o desenvolvimento de uma solução com risco tecnológico ainda não disponível no mercado são, respectivamente:
\begin{alternativas}
  \item $\reais{306.000,00}$ e Pregão Eletrônico por Menor Preço.
  \item $\reais{408.000,00}$ e Encomenda Tecnológica (ETEC, Lei n.\textsuperscript{o} 13.243/2016).
  \item $\reais{510.000,00}$ e Parceria Público-Privada de Concessão Rodoviária.
  \item $\reais{408.000,00}$ e Isenção do Simples Nacional.
  \item $\reais{612.000,00}$ e Encomenda Tecnológica (ETEC, Lei n.\textsuperscript{o} 13.243/2016).
\end{alternativas}

\textbf{Questão 5 (Cap. 4, Estado Empreendedor e Habitats de Inovação)}
Na análise crítica de políticas de CT\&I, a expressão ``catedral no deserto'' e a crítica de Mariana Mazzucato à ``socialização do risco com privatização dos ganhos'' referem-se, respectivamente, a:
\begin{alternativas}
  \item Parques tecnológicos integrados às cadeias locais e cobrança de \textit{royalties} públicos sobre patentes financiadas pelo Estado.
  \item Infraestrutura física de inovação desconectada da base produtiva local e apropriação privada integral de lucros gerados a partir de P\&D financiada com risco público.
  \item Incubadoras rurais autogestionárias e atuação de bancos comunitários solidários.
  \item Encomendas tecnológicas bem-sucedidas e participação acionária minoritária do setor público em \textit{startups}.
  \item Universidades públicas com NITs atuantes e políticas orientadas por missões climáticas.
\end{alternativas}

\textbf{Questão 6 (Cap. 5, Teoria da Mudança, SROI e \textit{Blended Finance})}
Um negócio de impacto em habitação sustentável captou $\reais{2.000.000,00}$ por meio de \textit{Blended Finance}: $25\%$ de capital filantrópico a $0\%$ a.a., $50\%$ de fundo \textit{impact-first} a $4\%$ a.a. e $25\%$ de banco comercial a $12\%$ a.a. O projeto gerou benefícios brutos de $\reais{4.000.000,00}$, com peso morto (\textit{deadweight}) estimado em $25\%$ e deslocamento nulo ($0\%$). Os valores corretos do custo ponderado de capital ($k_{blend}$) e do $SROI_{razao}$ líquido ajustado são:
\begin{alternativas}
  \item $k_{blend} = 4{,}0\%$ a.a. e $SROI_{razao} = 1{,}50$.
  \item $k_{blend} = 5{,}0\%$ a.a. e $SROI_{razao} = 1{,}50$.
  \item $k_{blend} = 5{,}0\%$ a.a. e $SROI_{razao} = 2{,}00$.
  \item $k_{blend} = 6{,}0\%$ a.a. e $SROI_{razao} = 1{,}25$.
  \item $k_{blend} = 5{,}3\%$ a.a. e $SROI_{razao} = 1{,}75$.
\end{alternativas}

\textbf{Questão 7 (Cap. 5, Quadro PCDO e \textit{Mission Drift})}
Sobre o Quadro PCDO (Austin et al., 2006) e os riscos do hibridismo organizacional em Negócios de Impacto, assinale a alternativa \textbf{correta}:
\begin{alternativas}
  \item O fenômeno de \textit{mission drift} ocorre quando a empresa social prioriza receitas de clientes de alta renda em detrimento do público vulnerável previsto em sua missão original.
  \item No Quadro PCDO, o \textit{Deal} (Negócio) é mais simples no empreendedorismo social porque o beneficiário sempre paga integralmente pelo serviço.
  \item Contratos de Impacto Social (SIBs) são imunes ao risco de seleção adversa (\textit{cream-skimming}).
  \item Negócios de impacto não precisam de modelo de receita sustentável, pois dependem exclusivamente de doações contínuas.
  \item A Teoria da Mudança considera \textit{outputs} e \textit{outcomes} como sinônimos exatos de faturamento bruto.
\end{alternativas}

\textbf{Questão 8 (Cap. 6, Tecnologia Social, IGs e Retenção de Valor Local)}
Uma cooperativa de cacau fino no Sul da Bahia comercializa amêndoas com Indicação de Procedência (IP) a $\reais{30,00}/\text{kg}$, frente ao preço de $\reais{20,00}/\text{kg}$ da \textit{commodity} comum. O custo fixo anual de rastreabilidade e certificação é de $\reais{40.000,00}$ e o volume comercializado é de $10.000\text{ kg/ano}$. O prêmio percentual de preço ($\pi_{IG}$), o ganho líquido unitário ($G_{unit}$) e o volume mínimo de equilíbrio ($Q^*$) são, respectivamente:
\begin{alternativas}
  \item $\pi_{IG} = 33{,}3\%$, $G_{unit} = \reais{6,00}/\text{kg}$ e $Q^* = 4.000\text{ kg/ano}$.
  \item $\pi_{IG} = 50{,}0\%$, $G_{unit} = \reais{10,00}/\text{kg}$ e $Q^* = 4.000\text{ kg/ano}$.
  \item $\pi_{IG} = 50{,}0\%$, $G_{unit} = \reais{6,00}/\text{kg}$ e $Q^* = 4.000\text{ kg/ano}$.
  \item $\pi_{IG} = 50{,}0\%$, $G_{unit} = \reais{6,00}/\text{kg}$ e $Q^* = 2.000\text{ kg/ano}$.
  \item $\pi_{IG} = 40{,}0\%$, $G_{unit} = \reais{4,00}/\text{kg}$ e $Q^* = 5.000\text{ kg/ano}$.
\end{alternativas}

\etapa{2}{Questões Dissertativas Críticas (Preparatório AV1)}

\textbf{Questão Dissertativa 1 (Governança, Estado Empreendedor e Território, Caps. 3 e 4)}
Um governo estadual pretende implantar um polo de bioeconomia e energias renováveis em uma região semiárida com forte presença de agricultura familiar, mas baixa densidade industrial. Explique por que a simples construção de um parque tecnológico padrão (isomorfismo institucional) tende a fracassar e proponha uma arquitetura de governança baseada na \textbf{Hélice Quíntupla} e em \textbf{Políticas Orientadas por Missões}, indicando como evitar a socialização do risco com privatização integral dos ganhos.
\espacoResolucao[3.0cm]

\textbf{Questão Dissertativa 2 (Impacto, Adequação Sociotécnica e Retenção de Valor, Caps. 5 e 6)}
Uma cooperativa extrativista analisa duas propostas de parceria para beneficiar frutos nativos da Caatinga: (A) instalação de uma fábrica automatizada de grande porte financiada por um fundo \textit{finance-first} sem trava de missão, comprando matéria-prima bruta dos associados; ou (B) adoção de unidades modulares de Tecnologia Social sob gestão cooperativa, com capital em \textit{Blended Finance} e busca de registro de Indicação Geográfica. Avalie comparativamente as duas propostas utilizando os conceitos de \textbf{Adequação Sociotécnica}, \textbf{\textit{Mission Drift}}, \textbf{Teoria da Mudança} e \textbf{Retenção de Valor Local ($RVL$)}.
\espacoResolucao[3.0cm]

\clearpage
\etapa{3}{Gabarito Comentado, Espelho Dissertativo e Diagnóstico de Revisão}

\begin{boxGabarito}[Gabarito Comentado das Questões Objetivas (1 a 8)]
\textbf{1. (B)} Trajetórias tecnológicas (Dosi) são cumulativas e dependem da densidade territorial de infraestrutura, conhecimento e normas descrita por Milton Santos.

\textbf{2. (B)} A eficiência coletiva combina ganhos espontâneos da aglomeração (externalidades marshallianas passivas) com ações conjuntas deliberadas entre firmas e instituições de apoio.

\textbf{3. (A)} Boschma alerta que proximidade excessiva (especialmente cognitiva e institucional) gera fechamento a novas ideias e aprisionamento (\textit{lock-in}) em trajetórias obsoletas.

\textbf{4. (B)} $E_{LeiBem} = 1.500.000 \times 0{,}80 \times 0{,}34 = \reais{408.000,00}$. Para contratação pública de solução com risco tecnológico inexistente no mercado, o instrumento adequado é a Encomenda Tecnológica (ETEC).

\textbf{5. (B)} ``Catedral no deserto'' designa infraestrutura física desconectada do território; a crítica de Mazzucato aponta a assimetria em que o Estado assume o risco inicial e o setor privado apropria integralmente os lucros sem contrapartida.

\textbf{6. (B)} $k_{blend} = (0{,}25 \times 0\%) + (0{,}50 \times 4\%) + (0{,}25 \times 12\%) = 0\% + 2{,}0\% + 3{,}0\% = 5{,}0\%$ a.a. Benefício líquido ajustado: $V_{liq} = 4.000.000 \times (1 - 0{,}25) \times 1{,}00 = \reais{3.000.000,00}$. Logo, $SROI_{razao} = 3.000.000 / 2.000.000 = 1{,}50$.

\textbf{7. (A)} O \textit{mission drift} ocorre quando pressões por rentabilidade deslocam o atendimento do público vulnerável para segmentos de maior renda.

\textbf{8. (C)} $\pi_{IG} = [(30 - 20)/20] \times 100\% = 50{,}0\%$. Custo unitário de certificação $= 40.000 / 10.000 = \reais{4,00}/\text{kg}$. $G_{unit} = (30 - 20) - 4{,}00 = \reais{6,00}/\text{kg}$. $Q^* = 40.000 / (30 - 20) = 4.000\text{ kg/ano}$.
\end{boxGabarito}

\begin{boxGabarito}[Espelho Dissertativo e Critérios de Pontuação (AV1)]
\textbf{Espelho da Questão Dissertativa 1 (5,0 pontos):}
\begin{itemize}[leftmargin=1.4em, itemsep=2pt]
  \item \textbf{Crítica ao isomorfismo e ``catedral no deserto'' (1,5 pt):} demonstrar que copiar infraestrutura de polos metropolitanos sem densidade empresarial e demanda local gera ociosidade e fuga de empresas.
  \item \textbf{Arquitetura em Hélice Quíntupla e Missão (2,0 pts):} mapear Universidade/IFs, empresas locais/cooperativas, Estado (compras públicas, ETEC, fomento), Sociedade Civil (agricultores familiares, H4) e Transição Socioecológica (convivência com o Semiárido e energia limpa, H5), definindo uma missão mensurável.
  \item \textbf{Mecanismos de retorno público (1,5 pt):} propor cláusulas de \textit{royalties} contingentes ao sucesso, co-titularidade de PI via NITs, participação acionária pública ou contrapartidas de permanência produtiva no território.
\end{itemize}

\textbf{Espelho da Questão Dissertativa 2 (5,0 pontos):}
\begin{itemize}[leftmargin=1.4em, itemsep=2pt]
  \item \textbf{Adequação Sociotécnica (1,5 pt):} demonstrar que a Proposta A reproduz Tecnologia Convencional incompatível com a escala e a autogestão local, enquanto a Proposta B aplica AST (módulos adaptados e manutenção local).
  \item \textbf{\textit{Mission Drift} e Financiamento (1,5 pt):} explicar que o capital \textit{finance-first} sem trava estatutária na Proposta A tende a comprimir o preço pago aos extrativistas; na Proposta B, o \textit{Blended Finance} reduz $k_{blend}$ sem exigir retorno especulativo.
  \item \textbf{Teoria da Mudança e $RVL$ (2,0 pts):} contrastar a venda de matéria-prima bruta (baixa $RVL$ na Proposta A) com a agroindustrialização cooperativa e proteção coletiva via IG (elevada $RVL$ e geração de \textit{outcomes} duradouros na Proposta B).
\end{itemize}
\end{boxGabarito}

\subsection*{Tabela de Diagnóstico de Revisão: Unidade I}

Utilize o quadro abaixo para mapear eventuais erros no simulado e direcionar sua revisão final antes da \textbf{AV1 (30\%)}:

\renewcommand{\arraystretch}{1.18}
\small
\begin{tabularx}{\textwidth}{@{}c p{4.2cm} X p{3.0cm}@{}}
  \toprule
  \rowcolor{etIndigo}
  \textcolor{white}{\bfseries Questão} & \textcolor{white}{\bfseries Capítulo de Referência} & \textcolor{white}{\bfseries Conceito-Chave Avaliado} & \textcolor{white}{\bfseries Meu Status} \\
  \midrule
  Q1 & Cap. 1, Base Tecnológica & Schumpeter, Dosi e Meio Técnico-Científico & $\square$ Dominei \quad $\square$ Revisar \\
  \rowcolor{bgGray}
  Q2 & Cap. 2, Aglomerações & Eficiência Coletiva em APLs & $\square$ Dominei \quad $\square$ Revisar \\
  Q3 & Cap. 3, Ecossistemas & Proximidades de Boschma e \textit{Lock-in} & $\square$ Dominei \quad $\square$ Revisar \\
  \rowcolor{bgGray}
  Q4 & Cap. 4, Governança CT\&I & Lei do Bem e Encomenda Tecnológica (ETEC) & $\square$ Dominei \quad $\square$ Revisar \\
  Q5 & Cap. 4, Estado Empreendedor & Catedral no Deserto e Socialização do Risco & $\square$ Dominei \quad $\square$ Revisar \\
  \rowcolor{bgGray}
  Q6 & Cap. 5, Negócios de Impacto & \textit{Blended Finance} ($k_{blend}$) e SROI Ajustado & $\square$ Dominei \quad $\square$ Revisar \\
  Q7 & Cap. 5, Quadro PCDO & \textit{Mission Drift} e Complexidade do \textit{Deal} & $\square$ Dominei \quad $\square$ Revisar \\
  \rowcolor{bgGray}
  Q8 & Cap. 6, IGs e Bioeconomia & Prêmio ($\pi_{IG}$), $G_{unit}$ e Ponto de Equilíbrio ($Q^*$) & $\square$ Dominei \quad $\square$ Revisar \\
  D1 & Caps. 3 e 4 (Integrada) & Hélice Quíntupla, Missões e Retorno Público & $\square$ Dominei \quad $\square$ Revisar \\
  \rowcolor{bgGray}
  D2 & Caps. 5 e 6 (Integrada) & AST, \textit{Mission Drift}, Teoria da Mudança e $RVL$ & $\square$ Dominei \quad $\square$ Revisar \\
  \bottomrule
\end{tabularx}
\normalsize
